% =====================================================================
%  2bat-ud5-14.tex  —  SESSIÓ 14: Estacions de probabilitat i distribucions
%  (sense preàmbul: s'inclou des de main.tex)
% =====================================================================
\horatitol{Sessió 14 — Estacions de probabilitat i distribucions}%
{Repàs actiu per quatre estacions, una per bloc de la unitat, i mapa final que connecta condicionada, binomial i normal.}

\subsection{Idea clau}

Repassem \textbf{tota la unitat} de manera activa, rotant per \textbf{quatre estacions},
una per cada bloc treballat. Al final construirem un \textbf{mapa de la unitat} que
connecta les eines —\textbf{condicionada}, \textbf{binomial} i \textbf{normal}— amb el
tipus de pregunta que resol cadascuna.

\begin{idea}
\textbf{Les quatre estacions}
\begin{enumerate}[leftmargin=1.6em, itemsep=2pt]
  \item \textbf{Probabilitat condicionada:} llegir una taula i calcular $P(A\,|\,B)$ amb
        la base correcta.
  \item \textbf{Binomial:} calcular $P(k)$ i el nombre esperat $\mu=n\cdot p$.
  \item \textbf{Normal — tipificar:} passar de $N(\mu,\sigma)$ a $Z$ i llegir la taula.
  \item \textbf{Normal — intervals:} calcular $P(a<X<b)$ i interpretar-ho.
\end{enumerate}
\end{idea}

\subsection{Les estacions}

\begin{exemple}[Estació 2 — binomial]
Cada usuari té un $p=0,5$ d'èxit; grup de $n=10$. La probabilitat que en tinguin
\textbf{exactament 6} és $P(6)=\binom{10}{6}(0,5)^6(0,5)^4=210\cdot 0,000977\approx
0,205=20,5\%$. El nombre esperat és $\mu=10\cdot 0,5=5$.
\end{exemple}

\begin{exemple}[Estació 3 — tipificar]
En una $N(100,15)$, tipifiquem $x=130$: $Z=\dfrac{130-100}{15}=2$. A la taula,
$P(Z<2)=0,9772$, així que $P(X<130)=97,72\%$.
\end{exemple}

\subsection{Exercicis resolts}

\begin{exresolt}[model de l'Estació 1: condicionada]
En una taula, $70$ usuaris viuen sols i, d'aquests, $40$ han trobat feina. Calcula
$P(\text{feina}\,|\,\text{viu sol})$.
\solucio
La base és el grup «viu sol» ($70$); dels quals $40$ tenen feina:
\[
P(\text{feina}\,|\,\text{viu sol})=\frac{40}{70}=\frac{4}{7}\approx 57,1\%.
\]
\resposta{$\dfrac{4}{7}\approx 57,1\%$.}
\end{exresolt}

\begin{exresolt}[model de l'Estació 4: interval normal]
En una $N(50,10)$, calcula $P(40<X<60)$.
\solucio
$z_{40}=\dfrac{40-50}{10}=-1$ i $z_{60}=\dfrac{60-50}{10}=1$.
\[
P(40<X<60)=P(Z<1)-P(Z<-1)=0,8413-0,1587=0,6826\approx 68,3\%.
\]
\resposta{$\approx 68,3\%$ (la regla del $68\%$).}
\end{exresolt}

\subsection{Exercicis per fer}

\noindent\textit{Una tasca per estació. Rota i resol-les totes.}

\begin{exercicis}
\begin{exlist}
  \item \textbf{(Estació 1)} En una taula, $40$ usuaris tenen més de 65 anys i,
        d'aquests, $25$ viuen sols. Calcula $P(\text{viu sol}\,|\,\text{més de 65})$.
  \item \textbf{(Estació 2)} Amb $n=8$ i $p=0,5$, calcula $P(\text{exactament }4)$.
        (Pista: $\binom{8}{4}=70$.)
  \item \textbf{(Estació 2)} Amb $n=20$ i $p=0,7$, calcula el nombre esperat d'èxits.
  \item \textbf{(Estació 3)} En una $N(60,10)$, tipifica $x=75$ i calcula $P(X<75)$.
        (Pista: $P(Z<1,5)=0,9332$.)
  \item \textbf{(Estació 4)} En una $N(20,4)$, calcula $P(16<X<24)$. Amb quina regla
        coincideix?
  \item \textbf{(Mapa de la unitat)} Completa el mapa connectant cada eina amb el tipus de
        pregunta que resol:
        \begin{center}
        \renewcommand{\arraystretch}{1.4}
        \begin{tabular}{p{4.0cm} p{8.4cm}}
        \toprule
        \textbf{Eina} & \textbf{Quina pregunta resol} \\
        \midrule
        Probabilitat condicionada & \rule{7.6cm}{0.4pt} \\
        Distribució binomial & \rule{7.6cm}{0.4pt} \\
        Distribució normal & \rule{7.6cm}{0.4pt} \\
        Tipificació ($Z$) & \rule{7.6cm}{0.4pt} \\
        \bottomrule
        \end{tabular}
        \end{center}
\end{exlist}
\end{exercicis}

% ---------------------------------------------------------------------
\fullsolucions{Sessió 14}
\begin{solucions}
\begin{sollist}
  \item Base: «més de 65» ($40$); viuen sols $25$.
        $P(\text{viu sol}\,|\,\text{més de 65})=\dfrac{25}{40}=\dfrac{5}{8}=0,625=62,5\%$.
  \item $P(4)=\binom{8}{4}(0,5)^4(0,5)^4=70\cdot 0,003906\approx 0,273=27,3\%$.
  \item $\mu=20\cdot 0,7=14$ èxits esperats.
  \item $Z=\dfrac{75-60}{10}=1,5$; $P(X<75)=P(Z<1,5)=0,9332=93,32\%$.
  \item $z_{16}=-1$, $z_{24}=1$; $P(16<X<24)=0,8413-0,1587=0,6826\approx 68,3\%$.
        Coincideix amb la regla del $68\%$.
  \item Resposta oberta. Idees: \textbf{condicionada} = probabilitat d'un cas sabent una
        condició (canviant la base); \textbf{binomial} = quants èxits en repeticions
        d'èxit/fracàs; \textbf{normal} = probabilitat d'una magnitud contínua (àrea sota
        la campana); \textbf{tipificació} = convertir qualsevol normal en $N(0,1)$ per
        usar una única taula.
\end{sollist}
\end{solucions}
